% GENERATED FILE. Edit canonical lessons, then run npm run generate:books.
% canonical-source-hash: fnv1a64:be5fe055e5f32735
% canonical-lessons: TA-C299-tippiti, TA-C299-parriyeri, TA-C299-tikkiraiyaku, TA-C299-tiyanaippu-virar, TA-C299-muccuttinaru

\chapter{Fire in the Night}
\label{ch:ta-fire-in-the-night}
\hlchaptermodality{299}

\begin{chapteropening}
\textbf{By the end of this chapter:} \emph{I can describe a fire and what it did.}

Complete the last lesson of chapter 299: I can describe a fire and what it did.
\end{chapteropening}

\section[\texorpdfstring{\ta{தீப்பிடி} (tīppiṭi)}{தீப்பிடி (tīppiṭi)}]{\ta{தீப்பிடி} (tīppiṭi) — to catch fire}

\label{lesson:TA-C299-tippiti}



\begin{quote}
Before the new one: say the Tamil for to escape alive, then the Tamil for first aid.
\end{quote}

\subsection*{You'll want to know: \ta{தீப்பிடி}}
\textbf{\ta{தீப்பிடி}} — \emph{tīppiṭi} — \textquotedblleft{}to catch fire\textquotedblright{}.

\begin{quote}
\ta{ஒருவேளை} \ta{வீடு} \ta{தீப்பிடித்தால்}, \ta{நாங்கள்} \ta{உடனே} \ta{வெளியே} \ta{ஓடுவோம்}.
\end{quote}

\begin{quote}
\emph{oruvēḷai vīṭu tīppiṭittāl, nāṅkaḷ uṭaṉē veḷiyē ōṭuvōm.}
\end{quote}

\begin{quote}
If the house were to catch fire, we would run out at once.
\end{quote}

Past (it …): \ta{தீப்பிடித்தது}, \emph{tīppiṭittatu}.

The thing on fire is the subject; also written as two words, \ta{தீப்} \ta{பிடி}.

\begin{cousinweb}
\textbf{\ta{தீப்பிடி}} (\emph{tīppiṭi}) is a Tamil compound built from \ta{தீ}, \textquotedblleft{}fire\textquotedblright{} (DEDR 3266), and \ta{பிடி}, \textquotedblleft{}to catch, seize\textquotedblright{} (DEDR 4148); the stop doubles after \ta{தீ}.
\end{cousinweb}

\begin{grammarlens}[title={What you've built}]
One more word for telling what happened.
\end{grammarlens}

\subsection*{Your turn}
\begin{itemize}
\raggedright
  \item \emph{Say it:} \emph{tīppiṭi}
  \item \emph{Say it:} \emph{tīppiṭi}, once more
  \item \emph{Say it:} \emph{oruvēḷai vīṭu tīppiṭittāl, nāṅkaḷ uṭaṉē veḷiyē ōṭuvōm.} — If the house were to catch fire, we would run out at once.
  \item \emph{Recall:} say the Tamil for to escape alive, then the Tamil for first aid, then say \emph{tīppiṭi} again
\end{itemize}

\begin{tcolorbox}[breakable,colback=teal!4,colframe=teal!35!black,title={Before you move on}]
What does \ta{தீப்பிடி} mean? (\textquotedblleft{}To catch fire\textquotedblright{}.) Say it once more.
\end{tcolorbox}
\section[\texorpdfstring{\ta{பற்றியெரி} (paṟṟiyeri)}{பற்றியெரி (paṟṟiyeri)}]{\ta{பற்றியெரி} (paṟṟiyeri) — to blaze}

\label{lesson:TA-C299-parriyeri}



\begin{quote}
Before the new one: say the Tamil for first aid, then the Tamil for to catch fire.
\end{quote}

\subsection*{You'll want to know: \ta{பற்றியெரி}}
\textbf{\ta{பற்றியெரி}} — \emph{paṟṟiyeri} — \textquotedblleft{}to blaze\textquotedblright{}.

\begin{quote}
\ta{கடை} \ta{இரவு} \ta{முழுவதும்} \ta{பற்றியெரிந்தது}.
\end{quote}

\begin{quote}
\emph{kaṭai iravu muḻuvatum paṟṟiyerintatu.}
\end{quote}

\begin{quote}
The shop blazed all night.
\end{quote}

Past (it …): \ta{பற்றியெரிந்தது}, \emph{paṟṟiyerintatu}.

From \ta{பற்றி}, having caught, and \ta{எரி}, to burn.

\begin{cousinweb}
\textbf{\ta{பற்றியெரி}} (\emph{paṟṟiyeri}) is a Tamil compound built from \ta{பற்று}, \textquotedblleft{}to catch hold, be kindled\textquotedblright{} (DEDR 4034), and \ta{எரி}, \textquotedblleft{}to burn, blaze\textquotedblright{} (DEDR 811).
\end{cousinweb}

\begin{grammarlens}[title={What you've built}]
One more word for telling what happened.
\end{grammarlens}

\subsection*{Your turn}
\begin{itemize}
\raggedright
  \item \emph{Say it:} \emph{paṟṟiyeri}
  \item \emph{Say it:} \emph{paṟṟiyeri}, once more
  \item \emph{Say it:} \emph{kaṭai iravu muḻuvatum paṟṟiyerintatu.} — The shop blazed all night.
  \item \emph{Recall:} say the Tamil for first aid, then the Tamil for to catch fire, then say \emph{paṟṟiyeri} again
\end{itemize}

\begin{tcolorbox}[breakable,colback=teal!4,colframe=teal!35!black,title={Before you move on}]
What does \ta{பற்றியெரி} mean? (\textquotedblleft{}To blaze\textquotedblright{}.) Say it once more.
\end{tcolorbox}
\section[\texorpdfstring{\ta{தீக்கிரையாகு} (tīkkiraiyāku)}{தீக்கிரையாகு (tīkkiraiyāku)}]{\ta{தீக்கிரையாகு} (tīkkiraiyāku) — to burn down; to be gutted by fire}

\label{lesson:TA-C299-tikkiraiyaku}



\begin{quote}
Before the new one: say the Tamil for to catch fire, then the Tamil for to blaze.
\end{quote}

\subsection*{You'll want to know: \ta{தீக்கிரையாகு}}
\textbf{\ta{தீக்கிரையாகு}} — \emph{tīkkiraiyāku} — \textquotedblleft{}to burn down; to be gutted by fire\textquotedblright{}.

\begin{quote}
\ta{பழைய} \ta{வீடு} \ta{முழுவதும்} \ta{தீக்கிரையானது}.
\end{quote}

\begin{quote}
\emph{paḻaiya vīṭu muḻuvatum tīkkiraiyāṉatu.}
\end{quote}

\begin{quote}
The whole old house was gutted by fire.
\end{quote}

Past (it …): \ta{தீக்கிரையானது}, \emph{tīkkiraiyāṉatu}.

Literally to become food for the fire: \ta{தீ} + \ta{இரை}, prey.

\begin{cousinweb}
\textbf{\ta{தீக்கிரையாகு}} (\emph{tīkkiraiyāku}) is a Tamil compound built from \ta{தீ}, \textquotedblleft{}fire\textquotedblright{} (DEDR 3266), \ta{இரை}, \textquotedblleft{}food, prey\textquotedblright{} (DEDR 490), and \ta{ஆகு}, \textquotedblleft{}to become\textquotedblright{}: to become the fire's prey.
\end{cousinweb}

\begin{grammarlens}[title={What you've built}]
One more word for telling what happened.
\end{grammarlens}

\subsection*{Your turn}
\begin{itemize}
\raggedright
  \item \emph{Say it:} \emph{tīkkiraiyāku}
  \item \emph{Say it:} \emph{tīkkiraiyāku}, once more
  \item \emph{Say it:} \emph{paḻaiya vīṭu muḻuvatum tīkkiraiyāṉatu.} — The whole old house was gutted by fire.
  \item \emph{Recall:} say the Tamil for to catch fire, then the Tamil for to blaze, then say \emph{tīkkiraiyāku} again
\end{itemize}

\begin{tcolorbox}[breakable,colback=teal!4,colframe=teal!35!black,title={Before you move on}]
What does \ta{தீக்கிரையாகு} mean? (\textquotedblleft{}To burn down; to be gutted by fire\textquotedblright{}.) Say it once more.
\end{tcolorbox}
\section[\texorpdfstring{\ta{தீயணைப்பு} \ta{வீரர்} (tīyaṇaippu vīrar)}{தீயணைப்பு வீரர் (tīyaṇaippu vīrar)}]{\ta{தீயணைப்பு} \ta{வீரர்} (tīyaṇaippu vīrar) — a firefighter}

\label{lesson:TA-C299-tiyanaippu-virar}



\begin{quote}
Before the new one: say the Tamil for to blaze, then the Tamil for to burn down.
\end{quote}

\subsection*{You'll want to know: \ta{தீயணைப்பு} \ta{வீரர்}}
\textbf{\ta{தீயணைப்பு} \ta{வீரர்}} — \emph{tīyaṇaippu vīrar} — \textquotedblleft{}a firefighter\textquotedblright{}.

\begin{quote}
\ta{தீயணைப்பு} \ta{வீரர்கள்} \ta{விரைவாக} \ta{வந்தார்கள்}.
\end{quote}

\begin{quote}
\emph{tīyaṇaippu vīrarkaḷ viraivāka vantārkaḷ.}
\end{quote}

\begin{quote}
The firefighters came quickly.
\end{quote}

Plural \ta{தீயணைப்பு} \ta{வீரர்கள்}. \ta{அணைப்பு} is putting out, from \ta{அணை}.

\begin{cousinweb}
\textbf{\ta{தீயணைப்பு} \ta{வீரர்}} (\emph{tīyaṇaippu vīrar}) is a Tamil phrase built from \ta{தீ}, \textquotedblleft{}fire\textquotedblright{} (DEDR 3266), \ta{அணைப்பு}, \textquotedblleft{}putting out\textquotedblright{}, from \ta{அணை}, \textquotedblleft{}to extinguish\textquotedblright{} (DEDR 121), and \ta{வீரர்}, \textquotedblleft{}a brave man\textquotedblright{}.
\end{cousinweb}

\begin{grammarlens}[title={What you've built}]
One more word for telling what happened.
\end{grammarlens}

\subsection*{Your turn}
\begin{itemize}
\raggedright
  \item \emph{Say it:} \emph{tīyaṇaippu vīrar}
  \item \emph{Say it:} \emph{tīyaṇaippu vīrar}, once more
  \item \emph{Say it:} \emph{tīyaṇaippu vīrarkaḷ viraivāka vantārkaḷ.} — The firefighters came quickly.
  \item \emph{Recall:} say the Tamil for to blaze, then the Tamil for to burn down, then say \emph{tīyaṇaippu vīrar} again
\end{itemize}

\begin{tcolorbox}[breakable,colback=teal!4,colframe=teal!35!black,title={Before you move on}]
What does \ta{தீயணைப்பு} \ta{வீரர்} mean? (\textquotedblleft{}A firefighter\textquotedblright{}.) Say it once more.
\end{tcolorbox}
\section[\texorpdfstring{\ta{மூச்சுத்திணறு} (mūccuttiṇaṟu)}{மூச்சுத்திணறு (mūccuttiṇaṟu)}]{\ta{மூச்சுத்திணறு} (mūccuttiṇaṟu) — to suffocate; to choke}

\label{lesson:TA-C299-muccuttinaru}



\begin{quote}
Before the new one: say the Tamil for to burn down, then the Tamil for a firefighter.
\end{quote}

\subsection*{You'll want to know: \ta{மூச்சுத்திணறு}}
\textbf{\ta{மூச்சுத்திணறு}} — \emph{mūccuttiṇaṟu} — \textquotedblleft{}to suffocate; to choke\textquotedblright{}.

\begin{quote}
\ta{புகையில்} \ta{பல} \ta{பேர்} \ta{மூச்சுத்திணறினார்கள்}.
\end{quote}

\begin{quote}
\emph{pukaiyil pala pēr mūccuttiṇaṟiṉārkaḷ.}
\end{quote}

\begin{quote}
Many people choked in the smoke.
\end{quote}

Past (I …): \ta{மூச்சுத்திணறினேன்}, \emph{mūccuttiṇaṟiṉēṉ}.

The one choking is the subject: \ta{அவன்} \ta{மூச்சுத்திணறினான்}.

\begin{cousinweb}
\textbf{\ta{மூச்சுத்திணறு}} (\emph{mūccuttiṇaṟu}) is a Tamil compound built from \ta{மூச்சு}, \textquotedblleft{}breath\textquotedblright{} (DEDR 4886), and \ta{திணறு}, \textquotedblleft{}to struggle, be choked\textquotedblright{}, whose origin is not settled in the sources at hand.
\end{cousinweb}

\begin{grammarlens}[title={What you've built}]
One more word for telling what happened.
\end{grammarlens}

\subsection*{Your turn}
\begin{itemize}
\raggedright
  \item \emph{Say it:} \emph{mūccuttiṇaṟu}
  \item \emph{Say it:} \emph{mūccuttiṇaṟu}, once more
  \item \emph{Say it:} \emph{pukaiyil pala pēr mūccuttiṇaṟiṉārkaḷ.} — Many people choked in the smoke.
  \item \emph{Recall:} say the Tamil for to burn down, then the Tamil for a firefighter, then say \emph{mūccuttiṇaṟu} again
\end{itemize}

\begin{tcolorbox}[breakable,colback=teal!4,colframe=teal!35!black,title={Before you move on}]
What does \ta{மூச்சுத்திணறு} mean? (\textquotedblleft{}To suffocate; to choke\textquotedblright{}.) Say it once more.
\end{tcolorbox}
